\subsection{The actual ordinary centralizer image in the 630-state quotient}

\begin{result}{A.4.38}{Proved}{The ordinary centralizer has a proper image in the named finite centralizer}
\label{thm:sphere10-ordinary-centralizer-image}
Retain the actual quotient $\rho:M\to L$ of A.4.32,
the marked geometric $H$ of A.4.36, and the actual
ordinary centralizer $\mathcal N=C_M(H)$ of A.1.
Then the two source motions specified below, together
with $H$, generate $\mathcal N$.  Their exact finite
images give
\[
\begin{gathered}
 |\rho(\mathcal N)|=2^{14}3^{302},\qquad
 |C_L(\rho(H))|=2^{27}3^{510},\\
 [C_L(\rho(H)):\rho(\mathcal N)]=2^{13}3^{208}.
\end{gathered}                                                  \tag{OCM.1}
\]
These are the image and centralizer of the same specified
source map $\rho$ and the same marked subgroup.
In particular the proper inclusion in (OCM.1) is an
explicit finite-image comparison.  It does not rule out
the centralizer-recognition refinements already constructed.
\end{result}

\begin{proof}
Any label permutation commuting with $H$ preserves the
three two-point orbits separately, since their point
stabilizers are the three distinct order-two subgroups
of $H$.  It also preserves the regular four-point orbit.
Thus $C_M(H)$ is the actual colored centralizer used
in A.1, with the same quotient map there.
\par\medskip\noindent\textbf{Two source motions and their exact descent.}\quad
Keep the six points $0,\infty,\pm1,\pm i$ stationary,
and move the regular quartet as
$\{z,-z,1/z,-1/z\}$ along the following two paths:
\[
 \gamma_0(\theta)=2e^{i\theta},\qquad
 \gamma_1(\theta)=\frac54+\frac34e^{i\theta},
                 \qquad 0\leq\theta\leq\pi.
                                                               \tag{OCM.2}
\]
Their endpoints are respectively $-2$ and $1/2$.
They avoid $0,\infty,\pm1,\pm i$ throughout.
The actual quotient by $H$ is the map
\[
             u(z)=\frac{(z^2+1)^2}{4z^2},
       \qquad u_0=u(2)=\frac{25}{16}.
                                                               \tag{OCM.3}
\]
It has branch values $0,1,\infty$ at the three fixed
pairs.  The actual centralizer descent of A.1 identifies
$\mathcal N/H$ with the fundamental group of their
complement based at $u_0$.

The first image path is the positively oriented ellipse
\[
 u(\gamma_0(\theta))
   =\frac12+\frac{17}{16}\cos(2\theta)
                         +i\frac{15}{16}\sin(2\theta),
                                                               \tag{OCM.4}
\]
which surrounds both $0$ and $1$.
The second image is a Jordan loop surrounding $1$
and excluding $0$.  Here are exact checks of this
claim and of the needed based free generators.
The interior of $\gamma_1$ lies in the open first
quadrant.  Equality of two $u$ values forces
$z'\in\{z,-z,1/z,-1/z\}$, by factoring the difference
of (OCM.3).  The other three points lie outside that
quadrant.  Hence the image of the interior is injective;
only the two endpoints, which are reciprocal, are
identified.
The factorization
$u-1=(z-1)^2(z+1)^2/(4z^2)$ gives winding number one
about $1$: the argument of $z-1$ changes from zero to
$\pi$, while those of $z+1$ and $z$ have zero net
change in the right half-plane.
The factorization
$u=(z-i)^2(z+i)^2/(4z^2)$ gives winding number zero
about $0$, using the continuous arguments of $z\pm i$
in the right half-plane.

Moreover $1/2<|\gamma_1(\theta)|<2$ for interior
parameters.  The full inverse image of the ellipse
in (OCM.4) is the union of the circles $|z|=2$ and
$|z|=1/2$, by the same four-element fibre formula.
The second loop consequently meets the first only at
$u_0$.  At the point $z=4/5+3i/5$ of $\gamma_1$ its
image is $16/25$, inside the ellipse.  Thus the entire
second loop, except its basepoint, is inside the first.
The two nested based Jordan loops form a free basis
of the twice-punctured plane: their complement on the
sphere has three regions, each containing exactly one
of $0,1,\infty$, and retracting these punctured regions
to their boundary graph gives the rank-two graph.
This proves generation in the actual group
$\mathcal N/H$, beyond the winding-number calculation.

An isotopy on the quotient sphere following either
image loop and fixing the three branch points lifts,
starting at the identity, to an $H$-equivariant isotopy.
It fixes the six ramification points individually and
follows exactly the four paths in (OCM.2).
Its endpoint therefore centralizes $H$.
Call the two endpoint classes $g_0,g_1$.
They generate a subgroup mapping isomorphically to
the free group $\mathcal N/H$: their images are the
just-proved free basis, so the source map from the
abstract free group is injective as well as onto its
image.  This subgroup has trivial intersection with
$H$, also seen from its pointwise fixing the six
ramification labels.  Hence
\[
           \mathcal N=H\times\langle g_0,g_1\rangle,
                 \qquad \langle g_0,g_1\rangle\cong F_2.
                                                               \tag{OCM.5}
\]

\par\medskip\noindent\textbf{Literal words and exact crossing checks.}\quad
Use precisely the projection chart and marking braid
$B$ of A.4.36.  In its signed-generator list convention,
the physical motion words are
\[
\begin{aligned}
 T_0={}&[2,3,2,4,9,8,7,6,5,4,3,2,1,6,7,8,
                              -3,-9,-4,-5,-4,1],\\
 T_1={}&[2,3,2,-6,9,8,4,7,6,8,-9,-4],\\
       &g_0=B T_0 B^{-1},\qquad g_1=B T_1 B^{-1}.
\end{aligned}                                                   \tag{OCM.6}
\]
These are source words of lengths $74$ and $64$.
The same rational-crossing method as A.4.36 certifies
them, with the following explicit rational paths.
For $t=\tan(\theta/2)$, write $z=\alpha/\beta$;
the two choices are
\[
 (\alpha,\beta)
  =\bigl(2(1-t^2)+4it,\,1+t^2\bigr),\qquad
  \bigl(4+t^2+3it,\,2(1+t^2)\bigr).
                                                               \tag{OCM.7}
\]
The four homogeneous pairs are
$(\alpha,\beta),(-\alpha,\beta),(\beta,\alpha),(-\beta,\alpha)$.
Their sphere vectors have numerator
$(2\operatorname{Re}(\alpha\overline\beta),
  2\operatorname{Im}(\alpha\overline\beta),
  |\alpha|^2-|\beta|^2)$ and positive denominator
$|\alpha|^2+|\beta|^2$.
Inserting the fixed rational chart gives polynomial
numerators and denominators for each real and imaginary
coordinate.
The chart pole has original coordinate $2+3i$.
For the first motion it is on neither radius-two nor
radius-one-half circle; for the second, none of
$2+3i,-2-3i,(2+3i)^{-1},-(2+3i)^{-1}$
lies on the prescribed upper-right arc.
Thus every projected denominator is positive.
The program additionally checks its positive-root
count is zero.

For every pair of labels the program retains the exact
horizontal and vertical difference polynomials, all
positive rational isolating intervals, the derivative
and vertical sign bounds, and the successive adjacent
label positions.  Exact Sturm counts certify that the
$22$ and $12$ recorded crossings exhaust the paths.
Endpoint leading coefficients give the prescribed
endpoint coordinates and permutations.  Finally all
four commutators of (OCM.6) with the two actual $H$
words vanish in all $630$ coordinates.
This identifies the source motions and also checks
their finite implementation.

\par\medskip\noindent\textbf{The full finite centralizer.}\quad
Put $T=N_3^{630}=\ker(L\to E_4)$.
For a state orbit of $H$ choose a representative
$\omega$.  A fixed tuple in $T$ is determined by its
value at $\omega$, which must centralize the actual
weight images of $H_\omega$ in $N_3$.
Every other value is the conjugate prescribed by
an actual $H$ transporter.  The stabilizer condition
is exactly the independence of this choice.
The complete finite check gives the table
\[
\begin{array}{c|r|r}
 \text{state-orbit size}&\text{number of such orbits}
                       &\text{allowed factor order}\\ \hline
 4&144&27\\
 2&21&27\\
 2&3&3\\
 1&3&27\\
 1&3&3
\end{array}
\qquad
 C_T(H)\cong N_3^{168}\times C_3^6.
                                                               \tag{OCM.8}
\]
All order-three allowed factors are the actual center
of $N_3$.  Every allowed value and all its transported
coordinates are retained and verified, proving the
displayed product and order $3^{510}$.

The permutation centralizer of $H$ is the actual
$C_2^5$ used in A.4.29: its three first generators
flip the three fixed pairs, and its last two translate
the regular quartet.  In $V_4$ the invariant equation
matrix is the actual map
$v\mapsto((1-a)v,(1-b)v)$.
Its integral-mod-four elimination has $23$ unit
pivots, one residual two-pivot, and $11$ zero
directions.  Hence $|V_4^H|=2^{23}$.
For each of the $32$ colored permutations, the
centralizer lift equations are solved using the actual
section coordinates of A.4.36 and the full cocycle
of A.4.31.  Precisely the $16$ permutations satisfying
$r_A+r_B+r_C=0$ lift; an explicit pure coordinate
for every lift is retained.  Thus
$|C_{E_4}(H)|=2^{27}$.

The centralizer projection onto this group is
surjective.  Indeed a lift in $L$ of a centralizer
element conjugates $H$ to a section over the same
$H$ in $E_4$.  Its difference from the original
section is a crossed homomorphism into $T$.
Modulo $Z(T)$, average its four values and multiply
by $4^{-1}$ in $\mathbb F_3$.
The cocycle equation expresses its value at $h$ as
$v-hv$; conjugation by a lift of $-v$ removes it.
The remaining difference is central, and the same
average in $Z(T)$ removes that difference.
This constructs an element of $T$ correcting the
lift to a pointwise centralizer.  Consequently
\[
 1\longrightarrow C_T(H)\longrightarrow C_L(H)
            \longrightarrow C_{E_4}(H)\longrightarrow1
                                                               \tag{OCM.9}
\]
is exact.  Its order is the second value in (OCM.1).

\par\medskip\noindent\textbf{The actual ordinary image and its retained kernel.}\quad
The exact image of $\langle g_0,g_1\rangle$ in $E_4$
has order $4096=2^{12}$.  A breadth-first census
retains all its source representatives and both
generator transitions.  Each of its $8192$ Schreier
edges gives a literal source kernel word.
Evaluating those words in $L$ gives elements of $T$,
and their membership in the fixed product (OCM.8)
is verified coordinate by coordinate.

Choose actual generators $x,y,z=xyx^{-1}y^{-1}$ of $N_3$.
The program checks every one of its $27^2$ products
in the coordinates
\[
 (i,j,k)(i',j',k')
       =(i+i',j+j',k+k'-ji')\quad\text{over }\mathbb F_3.
                                                               \tag{OCM.10}
\]
The change of central coordinate $k\mapsto k+2ij$
turns this product into addition plus one half of
the alternating commutator.  In particular the
subgroup generated by specified elements is exactly
the linear span of their changed coordinates and
all pairwise brackets.  This follows directly by
multiplying two elements and then the inverse half
power of their commutator; it uses only the
invertibility of two in $\mathbb F_3$.

Compress a fixed tuple to its $174$ orbit
representatives.  The $168$ full factors give
$336$ noncentral coordinates, and all $174$ factors
give a central coordinate, for a total of $510$.
Exact row reduction of the $8192$ actual Schreier
vectors gives noncentral rank $198$ and central
rank $98$.  Adjoining all brackets of the retained
noncentral basis adds precisely six central
directions.  The generated kernel therefore has
order $3^{198+104}=3^{302}$.
The matrices, pivots and six bracket witnesses are
retained, and every pairwise bracket is checked
against the resulting span.
Schreier's kernel generation proves that this is
the whole kernel of the two-motion image over
its $E_4$ image, not merely a subgroup of that kernel.
Its order is consequently $2^{12}3^{302}$.

Both motions fix the first six labels pointwise,
whereas every nonidentity element of $H$ moves
some of those labels.  Their finite images
therefore intersect trivially.  They commute,
so (OCM.5) gives
$|\rho(\mathcal N)|=4\cdot2^{12}3^{302}$.
This proves the first value and the exact index
in (OCM.1).

The three complete programs are
\path{check_sphere10_centralizer_motion.py},
\path{check_sphere10_geometric_fixed_kernel.py}, and
\path{check_sphere10_ordinary_centralizer_image.py},
all in \path{work/ghost_code/}.
Their actual text outputs, JSON source words and
matrix archives are preserved.
The statement concerns the fixed output $L$.
The separately constructed stronger finite input
may impose the ordinary-image constraint upon
projection; the proper inclusion computed here
does not contradict such a refinement.
\end{proof}
